\documentclass[10pt, twocolumn]{article}

% ─── Codificación y lenguaje ───────────────────────────────────────────────────
\usepackage[utf8]{inputenc}
\usepackage[T1]{fontenc}
\usepackage[spanish]{babel}

% ─── Matemáticas ──────────────────────────────────────────────────────────────
\usepackage{amsmath, amssymb, amsthm}

% ─── Tipografía y geometría ────────────────────────────────────────────────────
\usepackage{lmodern}
\usepackage[margin=1.8cm, top=2cm, bottom=2cm]{geometry}

% ─── Tablas y figuras ─────────────────────────────────────────────────────────
\usepackage{booktabs}
\usepackage{tabularx}
\usepackage{array}
\usepackage{graphicx}
\usepackage{float}

% ─── Color y cajas de advertencia ─────────────────────────────────────────────
\usepackage{xcolor}
\usepackage{tcolorbox}
\tcbuselibrary{breakable, skins}

% ─── Hipervínculos ────────────────────────────────────────────────────────────
\usepackage[hidelinks]{hyperref}

% ─── Cajas de calibración / notas ─────────────────────────────────────────────
\newtcolorbox{notacal}{
  colback=yellow!10, colframe=orange!60!black,
  fonttitle=\bfseries\small, title=Nota de calibración,
  breakable, left=4pt, right=4pt, top=3pt, bottom=3pt
}

% ─── Marcado de IA ────────────────────────────────────────────────────────────
\newcommand{\ia}[1]{\textcolor{teal}{\textbf{\{ia\}}~#1}}

% ─── Metadatos ────────────────────────────────────────────────────────────────
\title{\textbf{Análisis Dinámico del \emph{Thermal Throttling}\\
       en el Procesamiento Local de Inteligencia Artificial}}

\author{%
  % Completa con los nombres e integrantes del equipo
  David  Alexander Salazar Villa \and Julian David Ramirez Rodriguez \and Juan Andres Gaviria 
}

\date{\today}

% ══════════════════════════════════════════════════════════════════════════════
\begin{document}
\maketitle

% ─── Responsabilidades ────────────────────────────────────────────────────────
\noindent\textbf{Responsabilidades del equipo:}
\begin{itemize}
  \item \textbf{David  Salazar } — [nose]
  \item \textbf{Julian Ramirez} — [lol]
  \item \textbf{Juan Gaviria} — [xd]
\end{itemize}

% ─── Resumen ──────────────────────────────────────────────────────────────────
\begin{abstract}
Se modela la dinámica del \emph{thermal throttling} en procesadores bajo carga
de inteligencia artificial como un sistema de tres ecuaciones diferenciales no
lineales acopladas que relacionan temperatura, potencia e inferencia. El término
de throttling cúbico introduce una no linealidad que impide la solución analítica
y divide el comportamiento del sistema en dos regímenes: uno de aceleración libre
y uno de protección térmica. Se implementa el método de Runge-Kutta de cuarto
orden (RK4) desde cero para integrar el sistema, junto con otros dos métodos de
referencia, y se realiza un análisis de error variando el tamaño de paso para
caracterizar la degradación de convergencia inducida por los gradientes extremos
en la frontera de throttling. Se simulan cuatro escenarios de intervención
---base, \emph{undervolting}, mejora de disipación y cambio de pasta
térmica--- y se evalúa en cada uno si el sistema converge a un equilibrio estable
con inferencia viable. Los resultados se animan para ilustrar las trayectorias en
el espacio de estados.
\end{abstract}

% ══════════════════════════════════════════════════════════════════════════════
\section{Introducción}

\subsection{Presentación del sistema}

El sistema escogido es la unidad de procesamiento de una computadora (CPU/NPU)
operando bajo cargas de trabajo generadas por la inferencia de Modelos de
Lenguaje Grande (LLM). Este sistema es un entorno dinámico donde la potencia
eléctrica, la temperatura del silicio y la velocidad de procesamiento de la IA
(medida en \emph{tokens} por segundo) se acoplan. A diferencia de otros sistemas
mecánicos, las decisiones se toman en microsegundos y de forma no lineal, ya que
el sistema presenta un comportamiento de dos regímenes: una fase de aceleración
suave y una fase de desaceleración extrema. Estas fases emergen de la arquitectura
del hardware, que dispone de un firmware protector que, al alcanzar sus límites
térmicos, disminuye o corta abruptamente la fuente de energía hacia el procesador
para evitar daños físicos.

\subsection{Relevancia}

Con los estándares actuales de hardware para IA local (requisito industrial de
$>40$~TOPS), un equipo ejecutando un LLM está siempre en tensión física gracias
a tres fuerzas en competencia:

\begin{itemize}
  \item El sistema operativo busca maximizar la productividad asignando la mayor
        potencia posible para generar \emph{tokens} con baja latencia.
  \item El procesador, por efecto Joule, incrementa su temperatura de forma
        proporcional a esa carga de trabajo.
  \item El firmware interviene como mecanismo de seguridad recortando la energía
        de forma drástica para prevenir daños irreversibles en la arquitectura del chip.
\end{itemize}

Definir un equilibrio entre estas tres fuerzas asegura la viabilidad operativa
de cualquier equipo dedicado a la inteligencia artificial.

\subsection{Fenómeno de estudio}

El \emph{thermal throttling} es el fenómeno central. En términos de sistemas se
define como una penalización impuesta por el hardware para desacelerar el flujo de
datos y reducir la generación de calor. En el contexto de la IA, este fenómeno
genera oscilaciones severas: cuando el equipo entra en throttling, el rendimiento
se desploma por debajo de los umbrales mínimos de usabilidad (frecuentemente por
debajo de 5~tokens/s), para luego intentar recuperarse una vez la temperatura
desciende, creando un ciclo de inestabilidad que compromete la integridad del
procesamiento. Modelar este fenómeno permite predecir el ``punto de fatiga'' del
hardware y diseñar intervenciones que optimicen la longevidad del equipo sin
sacrificar su capacidad de cómputo.

\subsubsection{Por qué no existe solución analítica}

El sistema \textbf{no admite solución analítica exacta} por dos razones
complementarias:

El firmware actúa como un interruptor condicional que solo se activa cuando la
temperatura sobrepasa un límite crítico, dividiendo el espacio de estados en dos
regímenes con dinámicas cualitativamente distintas: un \emph{régimen libre}, donde
el sistema opera según las demandas de rendimiento, y un \emph{régimen de
protección}, donde se introducen fuerzas de frenado para preservar el silicio.
El término de throttling se representa como $\max(0, T - T_{\text{crit}})^3$,
ya que el firmware no realiza ninguna acción sino hasta que la temperatura supera
el umbral crítico.

El obstáculo matemático fundamental es que el instante exacto en el que ocurre
este cambio, $t^*$ tal que $T(t^*) = T_{\text{crit}}$, es una variable desconocida
que depende exclusivamente de la trayectoria previa del sistema. Al no poder
predecir cuándo ni cuántas veces el sistema alternará entre estados, es imposible
formular una expresión analítica que unifique ambos regímenes. Además, el calor
generado, la energía consumida y el rendimiento computacional están fuertemente
acoplados: la alteración de una variable afecta instantáneamente a las demás, lo
que impide desacoplar las ecuaciones para encontrar una solución general.

\subsection{Objetivos}

El objetivo principal es modelar matemáticamente y simular mediante métodos
numéricos la dinámica térmica y de rendimiento de un procesador bajo carga de IA,
identificando las condiciones bajo las cuales el sistema alcanza un estado de
equilibrio estable.

Como objetivos específicos, se busca: implementar un algoritmo RK4 programado
desde cero; realizar un análisis comparativo de convergencia y error variando el
tamaño de paso $h$ y validando contra un solucionador de alto orden; y animar las
trayectorias resultantes para interpretar visualmente la sensibilidad del hardware
ante diferentes parámetros de refrigeración y configuración de software.

\subsection{Estructura del informe}

La \textbf{Sección~\ref{sec:metodologia}} desglosa el modelo matemático y justifica
las constantes físicas. La \textbf{Sección~\ref{sec:numerica}} describe la
metodología numérica, validación y herramientas. La
\textbf{Sección~\ref{sec:resultados}} presenta las simulaciones de los cuatro
escenarios de intervención junto al análisis de error. La
\textbf{Sección~\ref{sec:conclusiones}} sintetiza los hallazgos y reflexiona sobre
el trabajo en equipo.

% ══════════════════════════════════════════════════════════════════════════════
\section{Metodología}
\label{sec:metodologia}

\subsection{Modelo matemático: sistema de ecuaciones diferenciales}

El fenómeno se modela como un sistema de tres ecuaciones diferenciales no lineales
de primer orden acopladas. Las tres variables de estado son:

\begin{itemize}
  \item $T(t)$: temperatura del procesador en el instante $t$ [\textdegree C]
  \item $P(t)$: potencia eléctrica asignada al procesador [W]
  \item $I(t)$: velocidad de inferencia (tokens generados por segundo) [tokens/s]
\end{itemize}

Las tres están acopladas: la potencia determina cuánto se calienta el procesador y
cuán rápido procesa; la temperatura activa el throttling que recorta la potencia;
la potencia recortada reduce la inferencia, lo que impulsa al sistema operativo a
demandar más potencia.

% ─── Ecuación 1 ───────────────────────────────────────────────────────────────
\subsubsection*{Ecuación 1: Dinámica Térmica}

\begin{equation}
  \frac{dT}{dt} = k_1 P(t) - k_2 \bigl(T(t) - T_{\text{amb}}\bigr)
  \label{eq:termica}
\end{equation}

Esta ecuación describe la \textbf{evolución de la temperatura del procesador}.
Es una aplicación directa del Modelo de Capacitancia Térmica Concentrada
(\emph{Lumped Capacitance Model}), donde el procesador se trata como un nodo único
con una masa térmica $C_{\text{th}}$ y una resistencia térmica $R_{\text{th}}$
hacia el ambiente. El balance de energía canónico es
$C_{\text{th}} \frac{dT}{dt} = \dot{E}_{\text{in}} - \dot{E}_{\text{out}}$,
donde la energía entrante es la potencia eléctrica disipada como calor (Efecto
Joule) y la saliente sigue la Ley de Enfriamiento de Newton. Dividiendo entre
$C_{\text{th}}$ se obtiene la forma exacta de la ecuación, con
$k_1 \equiv 1/C_{\text{th}}$ y $k_2 \equiv 1/(R_{\text{th}} C_{\text{th}})$.

El término $k_1 P(t)$ es la \textbf{tasa de calentamiento}: la potencia consumida
se disipa parcialmente como calor. La constante $k_1$ [\textdegree C/J] cuantifica
cuántos grados produce cada julio entregado; depende de la arquitectura del chip
a través de $\kappa_0$, que agrupa la ineficiencia térmica de la litografía del
silicio.

El término $-k_2(T - T_{\text{amb}})$ representa la \textbf{disipación térmica}
hacia el ambiente. La constante $k_2$ [s$^{-1}$] cuantifica la eficiencia del
sistema de disipación: cuando $T = T_{\text{amb}}$ no hay pérdida neta; cuando
$T \gg T_{\text{amb}}$ la disipación es máxima. Depende físicamente de
$\Phi_{\text{aire}}$ y $H_d$.

\smallskip
\noindent\textit{Referencia: Incropera, F.\,P.\ \& DeWitt, D.\,P.\ ---
\emph{Fundamentos de transferencia de calor y masa}, cap.\ Conducción transitoria.}

% ─── Ecuación 2 ───────────────────────────────────────────────────────────────
\subsubsection*{Ecuación 2: Dinámica de Inferencia}

\begin{equation}
  \frac{dI}{dt} = \gamma P(t) - \delta I(t)
  \label{eq:inferencia}
\end{equation}

Esta ecuación describe la \textbf{evolución de la velocidad de inferencia}.
Su forma proviene de los modelos de fluidos en teoría de colas
(\emph{Fluid Queueing Model}), cuya ecuación general es
$\dot{x} = \lambda(t) - \mu x(t)$, donde $\lambda(t)$ es la tasa de inyección
de trabajo y $\mu x(t)$ la fricción proporcional a la carga actual. La variable
$I(t)$ actúa como la variable de flujo $x(t)$.

El término $\gamma P(t)$ representa la \textbf{aceleración computacional}: a mayor
potencia, mayor frecuencia de operación y mayor \emph{throughput} de cálculo
matricial, operación central en los LLMs. La constante $\gamma$ mide la eficiencia
con que la energía se convierte en cómputo útil e incorpora las Instrucciones Por
Ciclo (IPC) y el número de hilos lógicos disponibles.

El término $-\delta I(t)$ representa la \textbf{fricción computacional}: ningún
sistema acelera indefinidamente. Al aumentar la velocidad de inferencia, se satura
el ancho de banda de memoria RAM y los buses de la placa base. La inferencia de
LLMs es un proceso severamente limitado por el ancho de banda de memoria
(\emph{memory bound}), por lo que $\delta$ es inversamente proporcional a la
velocidad del enlace más lento en la ruta de datos.

\smallskip
\noindent\textit{Referencia: Kleinrock, L.\ --- \emph{Queueing Systems, Vol.\,II:
Computer Applications}, cap.\ Fluid Approximations.}

% ─── Ecuación 3 ───────────────────────────────────────────────────────────────
\subsubsection*{Ecuación 3: Control de Potencia con Throttling No Lineal}

\begin{equation}
  \frac{dP}{dt} = \alpha\bigl(I_{\text{obj}} - I(t)\bigr)
                 - \beta \max\!\bigl(0,\, T(t) - T_{\text{crit}}\bigr)^3
  \label{eq:potencia}
\end{equation}

Esta ecuación es la \textbf{lógica de control de potencia} del sistema operativo y
el firmware, y contiene el término no lineal que hace al sistema matemáticamente
irresoluble de forma analítica.

El término $\alpha(I_{\text{obj}} - I(t))$ representa la \textbf{demanda del
sistema operativo}. En la teoría de control automático, corresponde exactamente a
la acción de un controlador proporcional (componente P de un PID): la ganancia
$\alpha$ actúa sobre el error $e(t) = I_{\text{obj}} - I(t)$ para ajustar la
potencia. Un $\alpha$ alto corresponde al perfil ``Máximo rendimiento''; un $\alpha$
bajo, a ``Ahorro de batería''.

El término $-\beta\max(0, T - T_{\text{crit}})^3$ modela el \textbf{throttling del
firmware}. Mientras $T \leq T_{\text{crit}}$ el término es exactamente cero y el
firmware no interviene; al superar $T_{\text{crit}}$, crece con la tercera potencia
del exceso $\Delta T = T - T_{\text{crit}}$.

Para verificar que no introduce discontinuidades problemáticas en $T = T_{\text{crit}}$,
se evalúan los límites laterales de
$\phi(\Delta T) = \max(0, \Delta T)^3$ y sus dos primeras derivadas en $\Delta T = 0$:

\begin{equation*}
  \lim_{\Delta T \to 0^{\pm}} \phi(\Delta T) = 0, \quad
  \lim_{\Delta T \to 0^{\pm}} \phi'(\Delta T) = 0, \quad
  \lim_{\Delta T \to 0^{\pm}} \phi''(\Delta T) = 0
\end{equation*}

Los tres límites bilaterales coinciden, confirmando que $\phi$ es continua junto
con sus dos primeras derivadas en el umbral. Para $p = 2$, $\phi'' = 2$ por la
derecha y $0$ por la izquierda ---una discontinuidad de curvatura que el integrador
experimentaría como un salto brusco. El exponente $p = 3$ es el mínimo entero que
elimina esta discontinuidad.

La elección del exponente cúbico se justifica por tres razones técnicas:

\smallskip\noindent\textbf{1. Comportamiento físico del hardware.}
La documentación técnica de Intel establece que las acciones ante una condición
térmica crítica escalan según la severidad del evento, desde reducir el ancho de
banda y la frecuencia hasta apagar subsistemas completos. Brooks y Martonosi, en
\emph{Dynamic Thermal Management for High-Performance Microprocessors}, documentan
que mecanismos como el \emph{clock gating} global o la reducción dinámica de
voltaje (DVFS) imponen caídas de rendimiento asimétricas cuya severidad aumenta
polinómicamente frente a la magnitud de la violación térmica. La curva cúbica
abstrae esa cascada: tolerante a micro-fluctuaciones en la frontera, pero
implacable ante excesos sostenidos.

\smallskip\noindent\textbf{2. Continuidad geométrica en el cruce de régimen.}
Nocedal y Wright, en \emph{Numerical Optimization}, establecen que exponentes
superiores al cuadrático se justifican cuando el método numérico requiere derivadas
continuas de orden mayor para no degradar su convergencia local. RK4 construye su
aproximación interpolando pendientes en puntos intermedios del intervalo
$[t_n, t_{n+1}]$; si la curvatura cambia abruptamente en $T_{\text{crit}}$, la
extrapolación interna introduce un error de truncamiento local que contamina la
solución en esa vecindad.

\smallskip\noindent\textbf{3. Coherencia con la teoría de penalización exterior.}
Rao, en \emph{Engineering Optimization: Theory and Practice}, explicita que aunque
$p = 2$ es el valor habitual por simplicidad analítica, el uso de $p = 3$ o
$p = 4$ es completamente legítimo cuando la estabilidad del método numérico exige
evitar discontinuidades de curvatura en el límite de la restricción. Adoptar
$p = 3$ produce además una fuerza restauradora más agresiva ante violaciones
grandes, sin necesidad de inflar artificialmente $\beta$.

La constante $\beta$ [W\,s$^{-1}$\,\textdegree C$^{-3}$] cuantifica el nivel de
severidad del firmware ante el calor. Un fabricante conservador usará un $\beta$
alto; uno orientado al rendimiento, un $\beta$ bajo.

\smallskip
\noindent\textit{Referencias: Ogata, K.\ --- \emph{Ingeniería de control moderna};
Ames et al.\ --- \emph{Control Barrier Functions: Theory and Applications};
Nocedal \& Wright --- \emph{Numerical Optimization}; Rao --- \emph{Engineering
Optimization}; Brooks \& Martonosi --- \emph{Dynamic Thermal Management for
High-Performance Microprocessors}.}

% ─── Constantes ───────────────────────────────────────────────────────────────
\subsection{Definición de constantes}

Para que el modelo refleje la realidad del hardware y permita simular
intervenciones físicas, las seis constantes de las ecuaciones se derivan
agrupando especificaciones tangibles de los componentes.

\subsubsection{Variables base de hardware}

\begin{table}[H]
\centering
\caption{Variables de hardware del modelo}
\label{tab:hardware}
\small
\begin{tabular}{@{}lll@{}}
\toprule
Símbolo & Unidad & Descripción \\
\midrule
$N_c$            & ---    & Núcleos físicos del procesador \\
$N_h$            & ---    & Hilos lógicos (\emph{threads}) \\
$f_p$            & GHz    & Frecuencia de reloj \\
$C_p$            & MB     & Caché del procesador \\
$v_{\text{ram}}$ & MT/s   & Velocidad de la RAM \\
$M_{\text{ram}}$ & GB     & Capacidad total de RAM \\
$v_{\text{alm}}$ & MB/s   & Velocidad de almacenamiento (NVMe) \\
$\Phi_{\text{aire}}$ & CFM & Flujo de aire de los ventiladores \\
$H_d$            & W/\textdegree C & Conductancia térmica del disipador \\
$v_{\text{bus}}$ & GT/s   & Velocidad del bus PCIe \\
\bottomrule
\end{tabular}
\end{table}

\subsubsection{Factores de escala empíricos}

Cada constante global se expresa como el producto de un \textbf{factor de escala
empírico} (subíndice $0$) y una combinación de variables de hardware. Los factores
de escala agrupan fenómenos microscópicos no modelados individualmente (resistencia
del silicio, geometría del chasis, \emph{overhead} del kernel, etc.) y ajustan las
unidades para garantizar la consistencia dimensional del modelo.

\medskip
\noindent\textbf{Dinámica Térmica:}

\begin{itemize}
  \item $k_1 = \kappa_0 \cdot (f_p \cdot N_c)$ [\textdegree C\,s$^{-1}$\,W$^{-1}$].
        A mayor frecuencia y más núcleos activos, más transistores conmutan
        simultáneamente y mayor disipación de energía como calor (Efecto Joule).
        El factor $\kappa_0$ captura la resistencia térmica específica del proceso
        litográfico del silicio (5\,nm, 7\,nm, etc.) y qué fracción de la energía
        se convierte en calor. Equivale a $1/C_{\text{th}}$ del modelo de
        Incropera: $k_1 \equiv 1/C_{\text{th}}$.

  \item $k_2 = \eta_0 \cdot (\Phi_{\text{aire}} \cdot H_d)$ [s$^{-1}$].
        La capacidad de enfriamiento depende del flujo de aire
        $\Phi_{\text{aire}}$ [CFM] y la conductancia térmica del disipador $H_d$
        [W/\textdegree C]: mayor flujo y mejor disipador producen un $k_2$ más alto.
        El factor $\eta_0$ corrige la geometría real del chasis (obstrucciones en
        las rejillas, flujo no laminar, presión ambiental) y garantiza la
        consistencia dimensional. Equivale a $1/(R_{\text{th}} C_{\text{th}})$.
\end{itemize}

\noindent\textbf{Dinámica de Inferencia:}

\begin{itemize}
  \item $\gamma = \gamma_0 \cdot (N_h \cdot f_p)$ [tokens\,s$^{-2}$\,W$^{-1}$].
        Basado en la Ley de Hierro del Rendimiento (\emph{Iron Law}), el
        \emph{throughput} bruto de cómputo es proporcional al número de hilos
        lógicos y a la frecuencia del reloj. El factor $\gamma_0$ representa las
        Instrucciones Por Ciclo (IPC) específicas para cálculo matricial,
        tipicamente acelerado por la NPU.

  \item $\delta = \delta_0 \,/\, \min(v_{\text{ram}},\, v_{\text{bus}})$
        [s$^{-1}$]. Según el Modelo Roofline, la generación de \emph{tokens} en
        LLMs opera en el régimen \emph{memory bound}: el cuello de botella real es
        el ancho de banda por el que los pesos del modelo llegan al procesador.
        El decaimiento del rendimiento es inversamente proporcional a la velocidad
        del enlace más lento, ya sea la RAM (MT/s) o el bus PCIe (GT/s). El factor
        $\delta_0$ cuantifica el \emph{overhead} intrínseco del sistema operativo.
\end{itemize}

\noindent\textbf{Control de Potencia:}

\begin{itemize}
  \item $\alpha = \alpha_0 \cdot N_h/N_c$
        [W\,s$^{-1}$\,(token\,s$^{-1}$)$^{-1}$]. La ratio $N_h/N_c$ es el
        índice de \emph{hyperthreading}: un procesador con HT activo expone más
        capacidad lógica al SO, induciendo al gobernador de energía a ser más
        agresivo. El factor $\alpha_0$ representa la sensibilidad del perfil
        energético: alto en ``Máximo rendimiento'', bajo en ``Ahorro de batería''.

  \item $\beta = \beta_0 \cdot f_p$
        [W\,s$^{-1}$\,\textdegree C$^{-3}$]. Los procesadores a mayor frecuencia
        son más vulnerables al daño por migración de electrones, por lo que el
        firmware escala su respuesta de corte con $f_p$. El factor $\beta_0$
        captura la filosofía de protección del fabricante: valores altos,
        fabricantes conservadores; valores bajos, orientados al rendimiento.
\end{itemize}

% ─── Análisis dimensional ─────────────────────────────────────────────────────
\subsection{Análisis dimensional}

Para garantizar la viabilidad de la integración numérica, el modelo debe ser
dimensionalmente homogéneo. Las unidades de las variables de estado son:
$T(t)$ [\textdegree C], $I(t)$ [tokens/s], $P(t)$ [W].

\medskip
\noindent\textbf{Ecuación térmica} ($dT/dt$ debe resultar en \textdegree C/s):
\[
\underbrace{k_1}_{\left[\frac{^\circ\text{C}}{\text{W\,s}}\right]}
\underbrace{P}_{\left[\text{W}\right]}
= \left[\frac{^\circ\text{C}}{\text{s}}\right], \qquad
\underbrace{k_2}_{\left[\text{s}^{-1}\right]}
\underbrace{(T - T_{\text{amb}})}_{\left[^\circ\text{C}\right]}
= \left[\frac{^\circ\text{C}}{\text{s}}\right]
\]

\noindent\textbf{Ecuación de inferencia} ($dI/dt$ en tokens/s$^2$):
\[
\underbrace{\gamma}_{\left[\frac{\text{tokens}}{\text{W\,s}^2}\right]}
\underbrace{P}_{\left[\text{W}\right]}
= \left[\frac{\text{tokens}}{\text{s}^2}\right], \qquad
\underbrace{\delta}_{\left[\text{s}^{-1}\right]}
\underbrace{I}_{\left[\frac{\text{tokens}}{\text{s}}\right]}
= \left[\frac{\text{tokens}}{\text{s}^2}\right]
\]

\noindent\textbf{Ecuación de potencia} ($dP/dt$ en W/s):
\[
\underbrace{\alpha}_{\left[\frac{\text{W}}{\text{tokens}}\right]}
\underbrace{(I_{\text{obj}}-I)}_{\left[\frac{\text{tokens}}{\text{s}}\right]}
= \left[\frac{\text{W}}{\text{s}}\right], \qquad
\underbrace{\beta}_{\left[\frac{\text{W}}{\text{s}\cdot^\circ\text{C}^3}\right]}
\underbrace{\Delta T^3}_{\left[^\circ\text{C}^3\right]}
= \left[\frac{\text{W}}{\text{s}}\right]
\]

% ══════════════════════════════════════════════════════════════════════════════
\section{Metodología Numérica}
\label{sec:numerica}

\subsection{Implementación de RK4}

Para resolver el sistema vectorial $\dot{\mathbf{x}} = \mathbf{f}(\mathbf{x}, t)$
con $\mathbf{x} = (T, I, P)^\top$, se implementa RK4 desde cero:

\begin{align*}
  \mathbf{k}_1 &= \mathbf{f}(\mathbf{x}_n, t_n) \\
  \mathbf{k}_2 &= \mathbf{f}\!\left(\mathbf{x}_n + \tfrac{h}{2}\mathbf{k}_1,\;
                                    t_n + \tfrac{h}{2}\right) \\
  \mathbf{k}_3 &= \mathbf{f}\!\left(\mathbf{x}_n + \tfrac{h}{2}\mathbf{k}_2,\;
                                    t_n + \tfrac{h}{2}\right) \\
  \mathbf{k}_4 &= \mathbf{f}(\mathbf{x}_n + h\,\mathbf{k}_3,\; t_n + h) \\[4pt]
  \mathbf{x}_{n+1} &= \mathbf{x}_n
    + \frac{h}{6}(\mathbf{k}_1 + 2\mathbf{k}_2 + 2\mathbf{k}_3 + \mathbf{k}_4)
\end{align*}

\noindent\textbf{Condición de frontera programática sobre $P(t)$.}
La ODE permite teóricamente que $P(t)$ sea negativa bajo throttling severo, lo cual
carece de sentido físico (el procesador siempre consume una potencia mínima de
inactividad $P_{\min}$). Al finalizar cada paso de integración se aplica la
corrección:
\begin{equation}
  P_{n+1} \leftarrow \max\!\bigl(P_{\min},\; P_{n+1}\bigr), \quad P_{\min} = 5\;\text{W}
\end{equation}
Esta corrección \textbf{no modifica la ODE}; es una regla lógica aplicada en el
código que descarta resultados matemáticamente válidos pero físicamente imposibles,
garantizando que el modelo permanezca en un dominio admisible.

\subsection{Inestabilidad inducida por gradientes extremos}

RK4 es un integrador \emph{explícito}: en cada paso calcula cuatro pendientes
locales y las combina para extrapolar el estado al siguiente instante. Esta
estrategia funciona bien cuando las pendientes cambian suavemente.

El problema surge cuando se activa el throttling. El término
$-\beta\max(0, T - T_{\text{crit}})^3$ genera una derivada respecto a $T$ igual a
$-3\beta(T - T_{\text{crit}})^2$, que crece con el \emph{cuadrado} del exceso
térmico: una superación de $10\,\text{\textdegree C}$ produce una pendiente
pronunciada, y $20\,\text{\textdegree C}$ la cuadruplica. El gradiente
\textbf{explota en función del estado del sistema}.

Cuando $h$ es demasiado grande, RK4 evalúa pendientes en puntos ya dentro de la
zona de gradiente extremo, sobreestimando masivamente el corte de potencia. Esto
produce que $P_{n+1}$ caiga a valores absurdos en un solo paso, colapsando $I$ y
disparando $\dot{P}$ en sentido contrario. El resultado es una oscilación numérica
divergente que no refleja el comportamiento físico real: es un artefacto del
integrador, no del sistema.

\subsection{Análisis de error y convergencia}

Para cada $h \in \{1.0,\; 0.5,\; 0.1,\; 0.05,\; 0.01,\; 0.005\}$ segundos, se
calcula el error relativo global respecto a RK45:

\begin{equation}
  E_{\text{rel}}(h) = \max_{t \in [0, T_{\text{sim}}]}
  \frac{\|\mathbf{x}_{\text{RK4}}(t;\,h) - \mathbf{x}_{\text{RK45}}(t)\|}
       {\|\mathbf{x}_{\text{RK45}}(t)\|}
\end{equation}

Se espera convergencia de orden 4 en el régimen suave ($E \sim h^4$) y degradación
del orden en la vecindad de $T = T_{\text{crit}}$. Esta degradación es un resultado
de análisis legítimo y debe discutirse en el informe.

\subsection{Validación}

La validación del RK4 implementado se realiza comparando sus soluciones contra
\texttt{scipy.integrate.solve\_ivp} con método RK45 de paso adaptativo. La
implementación se considera correcta cuando $E_{\text{rel}}$ disminuye con el orden
esperado al reducir $h$ en el régimen sin throttling activo ($T < T_{\text{crit}}$).
Adicionalmente, se verifica que RK4 y RK45 produzcan los mismos puntos de
equilibrio $(T^*, I^*, P^*)$ y el mismo comportamiento cualitativo en cada uno de
los cuatro escenarios.

\subsection{Herramientas}

La implementación se desarrolla en Python~3.x con las siguientes bibliotecas:
\texttt{NumPy} (operaciones vectoriales), \texttt{SciPy} (\texttt{solve\_ivp}/RK45,
solucionador de referencia), \texttt{Matplotlib} (visualización estática) y
\texttt{Manim} (animaciones). El código de los métodos numéricos está programado
directamente por el equipo, sin recurrir a solucionadores externos.

\subsection{Proceso de programación, simulación y animación}

La implementación sigue un orden secuencial. Primero se define la función vectorial
$\mathbf{f}(\mathbf{x}, t)$ que agrupa las tres ecuaciones, aplicando la condición
de frontera sobre $P_{\min}$ al final de cada evaluación. Sobre esta función se
construye el integrador RK4, que recibe $\mathbf{x}_n$ y $h$ y devuelve
$\mathbf{x}_{n+1}$ usando la fórmula de cuatro pendientes ponderadas.

Una vez verificado RK4 contra RK45 en el escenario base, se implementan los dos
métodos adicionales requeridos siguiendo la misma interfaz, lo que permite
intercambiarlos directamente en el bucle de integración para el análisis
comparativo.

Las simulaciones de los cuatro escenarios se ejecutan con el mismo horizonte
temporal y condiciones iniciales, variando únicamente los parámetros de hardware
correspondientes. Los resultados se grafican en el espacio de fases $(T, I, P)$
y como series temporales individuales. Las animaciones con Manim muestran la
trayectoria en tiempo real, resaltando el instante de activación del throttling,
la entrada al régimen oscilatorio y la convergencia al equilibrio en los escenarios
exitosos.

\subsection{Uso de IA}

Durante el desarrollo del proyecto se emplearon herramientas de inteligencia
artificial en los siguientes contextos:

\begin{itemize}
  \item \ia{Consulta de conceptos teóricos sobre el Modelo Roofline, el Modelo de
        Capacitancia Concentrada y las funciones de barrera de control, para
        clarificar su aplicabilidad al modelo propuesto.}
  \item \ia{Identificación de referencias bibliográficas relevantes (Incropera,
        Kleinrock, Nocedal \& Wright, Rao, Brooks \& Martonosi) para respaldar las
        decisiones de modelado.}
  \item \ia{Revisión de errores en la verificación dimensional de las constantes del
        modelo tras un análisis propio previo.}
  \item \ia{Sugerencia de condiciones iniciales y rangos de parámetros adicionales
        donde el sistema pudiera exhibir comportamientos interesantes, explorada
        después de forma propia mediante simulación.}
\end{itemize}

No se empleó IA para: generar el código de los métodos numéricos, escribir el
contenido del informe o la presentación, producir imágenes o animaciones, ni tomar
decisiones de modelado o análisis.

% ══════════════════════════════════════════════════════════════════════════════
\section{Resultados y Discusión}
\label{sec:resultados}

\subsection{Condiciones iniciales y parámetros}

Las condiciones iniciales representan el estado del equipo en $t = 0$, justo antes
de lanzar el proceso de inferencia (equipo encendido pero en reposo):

\begin{table}[H]
\centering
\caption{Condiciones iniciales}
\label{tab:ci}
\small
\begin{tabular}{@{}lcl@{}}
\toprule
Variable & Valor & Significado físico \\
\midrule
$T(0)$ & 45\,\textdegree C & Temperatura de reposo (carga base del SO) \\
$P(0)$ & 15\,W             & Potencia base del SO en inactividad \\
$I(0)$ & 0\,tokens/s       & LLM aún no iniciado \\
\bottomrule
\end{tabular}
\end{table}

\begin{table}[H]
\centering
\caption{Parámetros fijos del entorno}
\label{tab:params}
\small
\begin{tabular}{@{}lcl@{}}
\toprule
Parámetro & Valor & Descripción \\
\midrule
$T_{\text{amb}}$  & 25\,\textdegree C  & Temperatura ambiente \\
$T_{\text{crit}}$ & 90\,\textdegree C  & Umbral de activación del throttling \\
$I_{\text{obj}}$  & 20\,tokens/s       & Velocidad objetivo del gobernador \\
$I_{\text{min}}$  & 5\,tokens/s        & Umbral mínimo de viabilidad operativa \\
$P_{\text{min}}$  & 5\,W               & Potencia mínima física del procesador \\
\bottomrule
\end{tabular}
\end{table}

Se usa como referencia un portátil de gama media-alta típico para IA local
(procesador de 8 núcleos a 3.5\,GHz):

\begin{table}[H]
\centering
\caption{Valores de hardware de referencia}
\label{tab:hw}
\small
\begin{tabular}{@{}llcc@{}}
\toprule
Variable & Símbolo & Valor & Unidad \\
\midrule
Núcleos físicos     & $N_c$            & 8    & --- \\
Hilos lógicos       & $N_h$            & 16   & --- \\
Frecuencia          & $f_p$            & 3.5  & GHz \\
Caché               & $C_p$            & 24   & MB \\
Velocidad RAM       & $v_{\text{ram}}$ & 4800 & MT/s \\
Capacidad RAM       & $M_{\text{ram}}$ & 32   & GB \\
Almacenamiento      & $v_{\text{alm}}$ & 3500 & MB/s \\
Flujo de aire       & $\Phi_{\text{aire}}$ & 2.5 & CFM \\
Eficiencia disipador& $H_d$            & 5.0  & W/\textdegree C \\
Velocidad bus       & $v_{\text{bus}}$ & 16   & GT/s \\
\bottomrule
\end{tabular}
\end{table}

\begin{table}[H]
\centering
\caption{Factores de escala empíricos}
\label{tab:escala}
\small
\begin{tabular}{@{}llcl@{}}
\toprule
Factor & Símbolo & Valor & Justificación \\
\midrule
Ineficiencia térmica  & $\kappa_0$ & $0.002$           & Proceso 7\,nm \\
Convección            & $\eta_0$   & $0.004$           & Chasis parcialmente obstruido \\
Eficiencia IPC        & $\gamma_0$ & $0.015$           & NPU habilitada \\
Latencia intrínseca   & $\delta_0$ & $1.2\times10^{8}$ & Overhead estándar del SO \\
Ganancia gobernador   & $\alpha_0$ & $0.8$             & Perfil Máximo Rendimiento \\
Pánico firmware       & $\beta_0$  & $0.0001$          & Protección moderada \\
\bottomrule
\end{tabular}
\end{table}

Con estos valores, los parámetros del modelo resultan:

\begin{align*}
  k_1 &= 0.002 \times (3.5 \times 8) = 0.056\;\text{\textdegree C\,s}^{-1}\text{W}^{-1}\\
  k_2 &= 0.004 \times (2.5 \times 5.0) = 0.050\;\text{s}^{-1}\\
  \gamma &= 0.015 \times (16 \times 3.5) \times \ln\!\left(\tfrac{24 \times 16}{1}\right)
           \approx 5.03\;\text{tokens\,s}^{-2}\text{W}^{-1}\\
  \delta &= \frac{1.2\times10^{8}}{4800 \times 32 \times 3500} \approx 0.223\;\text{s}^{-1}\\
  \alpha &= 0.8 \times \frac{16}{8} = 1.6\;\text{W\,s}^{-1}\text{(tokens/s)}^{-1}\\
  \beta  &= 0.0001 \times 3.5 = 3.5\times10^{-4}\;\text{W\,s}^{-1}\text{\textdegree C}^{-3}
\end{align*}

\subsection{Análisis de puntos de equilibrio}

\subsubsection{Equilibrio sin throttling ($T^* < T_{\text{crit}}$)}

Igualando las tres derivadas a cero:

\begin{align*}
  \frac{dT}{dt}=0 &\implies T^* = T_{\text{amb}} + \frac{k_1}{k_2} P^* \\
  \frac{dI}{dt}=0 &\implies I^* = \frac{\gamma}{\delta} P^* \\
  \frac{dP}{dt}=0 &\implies I^* = I_{\text{obj}}
\end{align*}

De la tercera se obtiene $I^* = I_{\text{obj}}$. Sustituyendo:

\begin{equation}
  P^* = \frac{\delta \cdot I_{\text{obj}}}{\gamma}, \qquad
  T^* = T_{\text{amb}} + \frac{k_1 \delta \cdot I_{\text{obj}}}{k_2 \gamma}
\end{equation}

\begin{notacal}
Con los valores de referencia: $P^* = (0.223 \times 20)/5.03 \approx 0.887\;\text{W}$.
Este valor es físicamente inconsistente (un portátil bajo carga debería estar en
15--45\,W), lo que indica que $\gamma$ requiere recalibración iterativa. Este
proceso de ajuste es parte de la metodología y debe documentarse.
\end{notacal}

La condición de existencia del equilibrio estable es:

\begin{equation}
  T^* < T_{\text{crit}} \iff \frac{k_1 \delta}{k_2 \gamma} < \frac{T_{\text{crit}} - T_{\text{amb}}}{I_{\text{obj}}}
  \label{eq:viabilidad}
\end{equation}

Esta desigualdad es la \textbf{condición de viabilidad} del equipo para sostener
la carga de IA sin throttling.

\subsubsection{Estabilidad por retroalimentación negativa}

Para verificar que $(T^*, I^*, P^*)$ es un atractor estable, se analiza la
dirección de las derivadas ante perturbaciones.

\smallskip\noindent\textbf{Perturbación térmica.} Si la temperatura sube a
$T^* + \varepsilon$ con $\varepsilon > 0$:

\begin{equation*}
  \frac{dT}{dt}\bigg|_{T^*+\varepsilon}
  = \underbrace{k_1 P^* - k_2(T^*-T_{\text{amb}})}_{=\,0} - k_2\varepsilon = -k_2\varepsilon < 0
\end{equation*}

La derivada es negativa: el sistema enfría activamente hacia $T^*$. El coeficiente
$k_2 > 0$ actúa como fuerza de amortiguación proporcional al desvío.

\smallskip\noindent\textbf{Perturbación de inferencia.} Si $I$ cae por debajo
de $I^*$, el término $\alpha(I_{\text{obj}} - I)$ se vuelve positivo, aumentando
$\dot{P}$, lo que empuja $\dot{I} = \gamma P - \delta I$ hacia valores positivos,
recuperando la velocidad. El término $-\delta I(t)$ actúa simultáneamente como
fricción que impide sobrepasar $I^*$.

La estabilidad requiere que se satisfaga la misma condición~(\ref{eq:viabilidad}).

\subsubsection{Régimen oscilatorio}

Cuando la condición de viabilidad no se cumple, el sistema entra en el ciclo:

\begin{enumerate}
  \item El SO aumenta la potencia buscando alcanzar $I_{\text{obj}}$.
  \item La temperatura sube y supera $T_{\text{crit}}$.
  \item El firmware activa el throttling cúbico, colapsando la potencia.
  \item La temperatura baja al perder la fuente de calor.
  \item El SO intenta recuperar la potencia y el ciclo reinicia.
\end{enumerate}

Este régimen hace al equipo inviable para cargas de IA: la inferencia oscila entre
valores altos (inestables térmicamente) y valores bajos (en recuperación
post-throttling), sin converger a un estado estable con $I \geq I_{\min}$.

\subsection{Cuatro escenarios de intervención}

Los escenarios modifican los parámetros del hardware para evaluar qué
intervenciones permiten mover el sistema desde el régimen oscilatorio hacia un
equilibrio estable.

\subsubsection*{Escenario 0 --- Base (control)}
Portátil estándar sin modificaciones. Parámetros de la Tabla~\ref{tab:hw}.
Puede o no tener equilibrio estable según la condición~(\ref{eq:viabilidad}).
Es el punto de referencia para medir variaciones porcentuales.

\subsubsection*{Escenario 1 --- Intervención de software (undervolting)}
\textbf{Modificación:} Reducción de $\alpha_0$ (gobernador menos agresivo) y
limitación de $P_{\max}$ a 35\,W.\\
\textbf{Efecto esperado:} El SO sube la potencia más lentamente, dando tiempo al
disipador para mantenerse por debajo de $T_{\text{crit}}$. Puede estabilizar el
sistema a costa de reducir $I^*$.\\
\textbf{Variable de comparación:} ¿Alcanza $I^* \geq I_{\min} = 5$\,tokens/s?

\subsubsection*{Escenario 2 --- Hardware ligero (base refrigerante / limpieza)}
\textbf{Modificación:} Aumento de $\Phi_{\text{aire}}$ $\rightarrow$ $k_2$ sube.\\
\textbf{Efecto esperado:} Mayor disipación térmica desplaza $T^*$ hacia abajo,
permitiendo el equilibrio sin sacrificar velocidad.\\
\textbf{Variable de comparación:} Variación porcentual de $T^*$ y de la frecuencia
de oscilación respecto al escenario base.

\subsubsection*{Escenario 3 --- Hardware profundo (pasta térmica de metal líquido)}
\textbf{Modificación:} Aumento de $H_d$ y ajuste de $\eta_0$
$\rightarrow$ $k_2$ sube sustancialmente.\\
\textbf{Efecto esperado:} Si el nuevo $k_2$ cumple la condición~(\ref{eq:viabilidad}),
el sistema converge al equilibrio estable con inferencia plena.\\
\textbf{Variable de comparación:} Tiempo de convergencia y máxima temperatura
alcanzada durante la transición.

\subsection{Comparación entre métodos y tamaños de paso}

% PENDIENTE: completar con resultados de ejecución del código.
% Incluir tabla de E_rel(h) para cada h y cada método.
% Incluir gráfica log-log de E vs h mostrando la pendiente de orden 4.
[Sección por completar con resultados numéricos]

% ══════════════════════════════════════════════════════════════════════════════
\section{Conclusiones}
\label{sec:conclusiones}

% PENDIENTE: completar tras obtener los resultados de simulación.
[Sección por completar]

% ══════════════════════════════════════════════════════════════════════════════
\section*{Repositorio}

% Completar con el enlace al repositorio del equipo.
Código disponible en: \url{https://github.com/...}

% ══════════════════════════════════════════════════════════════════════════════
\begin{thebibliography}{99}

\bibitem{incropera}
Incropera, F.\,P.\ \& DeWitt, D.\,P.\ (2002).
\emph{Fundamentos de transferencia de calor y masa}.
Wiley. Cap.\ Conducción transitoria / Sistemas de capacitancia concentrada.

\bibitem{kleinrock}
Kleinrock, L.\ (1976).
\emph{Queueing Systems, Volume II: Computer Applications}.
Wiley-Interscience. Cap.\ Fluid Approximations.

\bibitem{ogata}
Ogata, K.\ (2010).
\emph{Ingeniería de control moderna}.
Pearson. Cap.\ Acciones básicas de control.

\bibitem{ames}
Ames, A.\,D., Coogan, S., Egerstedt, M., Notomista, G., Sreenath, K.,
\& Tabuada, P.\ (2019).
Control Barrier Functions: Theory and Applications.
\emph{18th European Control Conference (ECC)}.

\bibitem{nocedal}
Nocedal, J.\ \& Wright, S.\,J.\ (2006).
\emph{Numerical Optimization} (2nd ed.).
Springer.

\bibitem{rao}
Rao, S.\,S.\ (2019).
\emph{Engineering Optimization: Theory and Practice} (5th ed.).
Wiley.

\bibitem{brooks}
Brooks, D.\ \& Martonosi, M.\ (2001).
Dynamic Thermal Management for High-Performance Microprocessors.
\emph{Proceedings of the 7th International Symposium on High-Performance
Computer Architecture (HPCA)}.

\bibitem{hennessy}
Hennessy, J.\,L.\ \& Patterson, D.\,A.\ (2019).
\emph{Computer Architecture: A Quantitative Approach} (6th ed.).
Morgan Kaufmann.

\bibitem{williams}
Williams, S., Waterman, A., \& Patterson, D.\ (2009).
Roofline: An Insightful Visual Performance Model for Multicore Architectures.
\emph{Communications of the ACM}, 52(4), 65--76.

\end{thebibliography}

\end{document}
